\label{inv:C3.7}
\StatementRecord{inv:C3.7}{proved}{}{the hypotheses, notation and definitions stated here; \cref{inv:A3.2}; \cref{inv:C1.2}; \cref{inv:C2.1}; \cref{inv:eq:C-eta-negative}}

The threefold $X$ admits no pluriclosed (SKT) metric and no
astheno-K\"ahler metric.  The geometric input is the value $h=1$
and the global Dolbeault class $q$ in \ref{inv:A3.2}, whose restriction to a
regular fibre is $\eta=u\wedge w+6\gamma\wedge\delta$.
More generally, it is enough that the real image of
$H_A^{1,1}(X;\R)\to H_A^{1,1}(F;\R)$ be contained in
$\R\eta$ for one regular fibre $F$.
\end{proposition}


\begin{proof}
\begin{enumerate}[label=\arabic*.,leftmargin=*]
\item We first compute the sign that matters.  Put
$y=\operatorname{Im}\tau>0$, $m=\operatorname{Im}\mu$, and
$\Delta=\operatorname{Im}\beta-6m^2/y<0$.
For the marked period rows
$\sigma_1=6\mu\gamma+\tau u+w$ and
$\sigma_2=\beta\gamma+\mu u+\delta$ we have
$\eta=u\wedge\sigma_1+6\gamma\wedge\sigma_2$.
This proves that the real class $\eta$ has type $(1,1)$.
The determinant of the four real period columns, with rows ordered
as $\operatorname{Re}\zeta_1,\operatorname{Im}\zeta_1,
\operatorname{Re}\zeta_2,\operatorname{Im}\zeta_2$, is
$y\Delta<0$.  Since $\eta^2=12\gamma\wedge u\wedge w\wedge\delta$,
\begin{equation}\label{inv:eq:C-eta-negative}
\int_{F,\mathrm{or}_{\C}}\eta^2=-12.
\end{equation}
Equivalently the invariant representative of $\eta$ has Hermitian
coefficient matrix
\[
\frac1{y\Delta}
\begin{pmatrix}\operatorname{Im}\beta&-6m\\-6m&6y\end{pmatrix},
\qquad \det=\frac6{y\Delta}<0,
\]
so it has one positive and one negative eigenvalue.  The exact
symbolic determinant and an exact rational example are checked in
\texttt{code/c\_sign\_and\_fibre.py} by both eigenvalues and
diagonal elimination.

\item The explicit invariant fibre representative of $\gamma^{0,1}$ is
\[
\Gamma=\frac{i}{2\Delta}
\left(d\bar\zeta_2-\frac myd\bar\zeta_1\right).
\]
Writing $Q_F$ for the invariant $(1,1)$ representative of $\eta$,
direct exterior multiplication gives
\[
Q_F\wedge\Gamma=
\frac1{4y\Delta}
d\zeta_1\wedge d\bar\zeta_1\wedge d\bar\zeta_2,
\quad
\int_FQ_F\wedge\Gamma\wedge d\zeta_2=-1.
\]
This also proves the nonvanishing used in \ref{inv:C1.2}.

\item When $h=1$, \ref{inv:C2.1} identifies $H_A^{1,1}(X)$ with the sum of the
one horizontal source and its vertical conjugate.  A representative
$Q$ of $q$ is nonzero in the horizontal source, since its fibre
restriction is nonzero whereas $d_1g$ restricts to zero.  Thus
$[Q]_A,[\bar Q]_A$ form a basis.  Their restrictions to $F$ both
equal the real class $\eta$.  The real image of the restriction map
is therefore $\R\eta$.

\item If $\omega$ were a pluriclosed positive real $(1,1)$ form, its
restriction to $F$ would be pluriclosed and positive.  Average this
restriction under translations of the compact complex torus $F$.
The average $\omega_{\mathrm{av}}$ is invariant, closed, and
strictly positive.  Translation preserves the Aeppli class: on a
torus the $\partial\bar\partial$-lemma identifies Aeppli
cohomology with the invariant $(p,q)$ forms, so the average is the
unique invariant representative of that class.  This identification
can also be checked Fourier mode by Fourier mode: each nonzero mode
is a Koszul square for two independent type components of its real
nonzero covector and contributes zero Aeppli cohomology.  Thus
$[\omega_{\mathrm{av}}]=r\eta$ for some real $r$.
But positivity gives $\int_F\omega_{\mathrm{av}}^2>0$, whereas
\eqref{inv:eq:C-eta-negative} gives $-12r^2\le0$, a contradiction.
In dimension three the astheno-K\"ahler equation
$\partial\bar\partial\omega^{n-2}=0$ is exactly the SKT equation.
\end{enumerate}
\end{proof}